\documentclass[aspectratio=169]{beamer}
\usepackage{amsmath, amssymb, bm}

\setbeamertemplate{footline}[frame number]

\title{Robô de Codificação de Bobina}
\author{Clara Temponi Marigo, \\
        Eduardo Henrique Basilio de Carvalho, \\
        João Vitor Braga da Silva Alves, \\
        Lucas Rocha de Almeida do Carmo, \\
        Luiz Miguel Santiago, \\
        Renan Neves da Silva
        }
\date{}

\begin{document}

\frame{\titlepage}

\begin{frame}{Instanciação da Superfície}
Uma superfície cilíndrica $C$ é instanciada com uma geometria delimitadora aleatória e uma base ortonormal.
\begin{equation}
    R \sim \mathcal{U}(0.1, 0.6), \quad W \sim \mathcal{U}(0.4, 1.0)
\end{equation}
\begin{equation}
    B_{cyl} = [u_{axis}, v_{tangent}, n_{ref}]
\end{equation}
onde $R$ é o raio do cilindro, $W$ é a largura do cilindro, $B_{cyl}$ é o sistema de referência espacial, $u_{axis}$ é o eixo longitudinal, $v_{tangent}$ é a tangente transversal e $n_{ref}$ é a normal de referência.
\end{frame}

\begin{frame}{Trajetória de Engajamento}
Uma transição inicial engaja o efetuador final da configuração inicial para a superfície usando um polinômio de jerk mínimo.
\begin{equation}
    p(\tau) = p_{init} + \sigma(\tau)(p_{home} - p_{init})
\end{equation}
onde $p_{init} \in \mathbb{R}^3$ é a pose inicial de escrita na superfície e $p_{home} \in \mathbb{R}^3$ é a configuração espacial inicial.
\end{frame}

\begin{frame}{Geração Contínua de Traços}
Para cada caractere na sequência de entrada, os pontos de cada traço são interpolados em uma curva suave para calcular as correspondentes coordenadas planares.
\begin{equation}
    S_{2D}(t) = \text{make\_interp\_spline}(P)
\end{equation}
\begin{equation}
    (u(s), v(s)) = S_{2D}(t(s))
\end{equation}
onde $P \in \mathbb{R}^{M \times 2}$ é a matriz de coordenadas do traço, $S_{2D}(t)$ é a spline cúbica contínua e $(u(s), v(s))$ são as coordenadas bidimensionais avaliadas no comprimento de arco $s$.
\end{frame}

\begin{frame}{Projeção na Superfície}
As coordenadas parametrizadas por arco são projetadas na superfície cilíndrica para avaliar a trajetória espacial tridimensional.
\begin{equation}
    \theta(s) = \frac{v(s) - H/2}{R}
\end{equation}
\begin{equation}
    p(s) = c_{cyl} + \left(u(s) - \frac{W}{2}\right)u_{axis} + R \left( \cos(\theta(s))n_{ref} + \sin(\theta(s))v_{tangent} \right)
\end{equation}
onde $\theta(s)$ é o deslocamento angular ao longo do eixo transversal, $H$ é a altura do cilindro, $c_{cyl}$ é o centro geométrico do cilindro e $p(s) \in \mathbb{R}^3$ é a posição cartesiana resultante.
\end{frame}

\begin{frame}{Transições de Jerk Mínimo}
Sequências de transição interpolam dentro do espaço de parâmetros da superfície através de três fases: retração ortogonal, translação e engajamento com a superfície.
\begin{equation}
    \sigma(\tau) = -2\tau^3 + 3\tau^2
\end{equation}
\begin{equation}
    \xi(\tau) = \xi_0 + \sigma(\tau)(\xi_f - \xi_0)
\end{equation}
onde $\tau \in [0, 1]$ é o tempo normalizado, $\sigma(\tau)$ é o perfil escalar de jerk mínimo e $\xi = [u, v, d]^T \in \mathbb{R}^3$ é o estado contínuo de parâmetros da superfície contendo as coordenadas de superfície e o deslocamento espacial ortogonal.
\end{frame}

\begin{frame}{Trajetória de Retração}
Uma transição terminal retrai o efetuador final da superfície de volta para uma configuração cartesiana inicial.
\begin{equation}
    p(\tau) = p_{end} + \sigma(\tau)(p_{home} - p_{end})
\end{equation}
onde $p_{end} \in \mathbb{R}^3$ é a pose final de escrita na superfície e $p_{home} \in \mathbb{R}^3$ é a configuração espacial retraída.
\end{frame}

\begin{frame}{Cinemática Direta}
A malha de controle em tempo discreto mapeia o vetor de configuração de juntas para os sistemas de coordenadas espaciais do efetuador final usando as funções \texttt{jac\_geo} e \texttt{fkm} do módulo \texttt{uaibot}.
\begin{equation}
    J_{g,k}, x_k, z_k = \text{FK}(q_k)
\end{equation}
onde $q_k \in \mathbb{R}^7$ é a configuração de juntas no passo $k$, $x_k \in \mathbb{R}^3$ é a posição cartesiana do efetuador final, $z_k \in \mathbb{S}^2$ é o vetor unitário do eixo z e $J_{g,k} \in \mathbb{R}^{6 \times 7}$ é o Jacobiano geométrico espacial.
\end{frame}

\begin{frame}{Erro de Tarefa e Jacobiano}
O erro de tarefa agrega a diferença de rastreamento translacional e a perda de orientação escalar.

\begin{equation}
    e = \begin{bmatrix} x_k - x_{des} \\ 1 - z_{des}^T z_k \end{bmatrix}
\end{equation}
\begin{equation}
    J_{r} = \begin{bmatrix} J_{g, 0:3} \\ z_{des}^T S(z_k) J_{g, 3:6} \end{bmatrix}
\end{equation}
onde $e \in \mathbb{R}^4$ é o erro de tarefa composto, $x_{des}$ e $z_{des}$ são a posição e orientação de referência, $J_r \in \mathbb{R}^{4 \times 7}$ é o Jacobiano de tarefa projetado e $S(\cdot)$ constrói a matriz anti-simétrica.
\end{frame}

\begin{frame}{Velocidade Restauradora de Junta}
Uma velocidade restauradora artificial é calculada para forçar a descida de gradiente linear, afastando a configuração do manipulador dos limites absolutos de hardware.
\begin{equation}
    \nabla H_i = \frac{q_{k,i} - \bar{q}_i}{(q_{max, i} - q_{min, i})^2}
\end{equation}
\begin{equation}
    \dot{q}_0 = -k_0 \nabla H
\end{equation}
onde $\nabla H \in \mathbb{R}^7$ é o gradiente de penalidade de limite do inverso do quadrado, $q_{min}$ e $q_{max}$ são os limites da junta, $\bar{q}$ é a configuração mediana, $k_0 \in \mathbb{R}$ é o ganho de centralização e $\dot{q}_0 \in \mathbb{R}^7$ é a velocidade restauradora da junta.
\end{frame}

\begin{frame}{Formulação do Objetivo de Otimização}
A função de custo penaliza tanto a folga no rastreamento da tarefa quanto os desvios de velocidade no espaço nulo.
\begin{equation}
    x_{opt} = \begin{bmatrix} \dot{q}^T & w^T \end{bmatrix}^T
\end{equation}
\begin{equation}
    J_{cost}(\dot{q}, w) = \frac{1}{2} \left( \lambda_d \|\dot{q}\|^2 + \|w\|^2 \right) - \lambda_d \dot{q}_0^T \dot{q}
\end{equation}
onde $x_{opt} \in \mathbb{R}^{11}$ é o vetor de decisão combinado, $\dot{q} \in \mathbb{R}^7$ são as velocidades das juntas, $w \in \mathbb{R}^4$ são as variáveis de folga do espaço de tarefa e $\lambda_d \in \mathbb{R}$ é o fator de amortecimento da velocidade da junta.
\end{frame}

\begin{frame}{Restrição de Igualdade da Tarefa}
O objetivo primário de rastreamento é formulado como uma igualdade linear restrita pelas variáveis de folga.
\begin{equation}
    J_r \dot{q} - w = -K e + v_{des}
\end{equation}
onde $K \in \mathbb{R}^{4 \times 4}$ é a matriz diagonal positivo-definida de rigidez da tarefa e $v_{des} \in \mathbb{R}^4$ é a velocidade "de pose" desejada de feedforward.
\end{frame}

\begin{frame}{Distância Espacial da Superfície}
A distância escalar ortogonal até o limite da superfície é avaliada para impor estritamente a prevenção de colisões.
\begin{equation}
    v_{rad} = (x_k - c_{cyl}) - \left( (x_k - c_{cyl})^T u_{axis} \right) u_{axis}
\end{equation}
\begin{equation}
    d = \|v_{rad}\| - R
\end{equation}
onde $v_{rad} \in \mathbb{R}^3$ é o vetor radial ortogonal ao eixo principal e $d \in \mathbb{R}$ é a distância mínima sinalizada para a superfície.
\end{frame}

\begin{frame}{Restrições de Desigualdade}
A segurança espacial e as restrições de hardware definem limites rígidos de desigualdade linear nas velocidades das juntas.
\begin{equation}
    -n(x_k)^T J_{g, 0:3} \dot{q} \leq \gamma (d - d_{safe})
\end{equation}
\begin{equation}
    \max \left(-\dot{q}_{max}, \frac{q_{min}-q_k}{\Delta t} + \epsilon \right) \leq \dot{q} \leq \min \left(\dot{q}_{max}, \frac{q_{max}-q_k}{\Delta t} - \epsilon \right)
\end{equation}
onde $n(x_k) \in \mathbb{S}^2$ é a normal externa à superfície, $\gamma$ é a taxa de convergência da barreira, $d_{safe}$ é a margem espacial estrita, $\dot{q}_{max}$ é o limite de velocidade do atuador, $\Delta t$ é o passo de integração e $\epsilon$ é uma tolerância numérica.
\end{frame}

\begin{frame}{Integração Cinemática}
A derivada de estado ideal é avaliada usando a função \texttt{solve\_qp} do módulo \texttt{qpsolvers} e integrada para atualizar a configuração do manipulador.
\begin{equation}
    \dot{q}^* = \text{solve\_qp}(J_{cost}, \text{constraints})
\end{equation}
\begin{equation}
    q_{k+1} = q_k + \dot{q}^* \Delta t
\end{equation}
onde $\dot{q}^* \in \mathbb{R}^7$ é a solução ótima de velocidade da junta e $q_{k+1} \in \mathbb{R}^7$ é a configuração sequencial atualizada da junta.
\end{frame}

\begin{frame}{Simulação}
As etapas desde o calculo da cinemática direta até a integração cinemática são repetidas em um loop de controle em tempo discreto até que a trajetória seja completamente rastreada.
A simulação é renderizada e salva pela biblioteca \texttt{uaibot} para visualização.
\end{frame}

\end{document}
